\chapter{Why the Search Abandons a Move}
\label{ch:refutation}

We examine how the search handles bad candidate moves across expanded visit budgets.

\section{The Refutation Ladder}

\begin{realdata}[Refutation Ladder]
\begin{center}
\begin{minipage}{\linewidth}
\captionsetup{type=figure}
\centering
\subcaptionbox{64 nodes: Kd6 N=31, Kf6 N=24}{% GENERATED by tools/make_figures.py — do not edit.
\begin{tikzpicture}
  \node[vgroot] (root) at (0,0) {Root\\ \scriptsize Budget: 64 nodes};
  \node[vgnode, fill=BuildBlue!45, inner sep=1.8pt] (kd6) at (-4.5,-2.2) {Kd6\\ \scriptsize $n{=}31$\ \ $Q{=}0.99344$};
  \node[vgnode, fill=BuildBlue!35, inner sep=1.8pt] (kf6) at (-1.5,-2.2) {Kf6\\ \scriptsize $n{=}24$\ \ $Q{=}0.95335$};
  \node[vgunvisited, inner sep=1.8pt] (kf5) at (1.5,-2.2) {Kf5\\ \scriptsize $n{=}0$\ \ $Q_{\mathrm{FPU}}{=}0.664$};
  \node[vgunvisited, inner sep=1.8pt] (kd5) at (4.5,-2.2) {Kd5\\ \scriptsize $n{=}0$\ \ $Q_{\mathrm{FPU}}{=}0.664$};
  \vgedgewidth{31}{31}\draw[vgedge, line width=\vgEW pt] (root) -- (kd6);
  \vgedgewidth{24}{31}\draw[vgedge, line width=\vgEW pt] (root) -- (kf6);
  \vgedgewidth{0}{31}\draw[vgedge, line width=\vgEW pt] (root) -- (kf5);
  \vgedgewidth{0}{31}\draw[vgedge, line width=\vgEW pt] (root) -- (kd5);
\end{tikzpicture}
}\\[2ex]
\subcaptionbox{128 nodes: Kf5 \& Kd5 visited once ($Q=0.000$)}{% GENERATED by tools/make_figures.py — do not edit.
\begin{tikzpicture}
  \node[vgroot] (root) at (0,0) {Root\\ \scriptsize Budget: 128 nodes};
  \node[vgnode, fill=BuildBlue!50, inner sep=1.8pt] (kd6) at (-4.5,-2.2) {Kd6\\ \scriptsize $n{=}62$\ \ $Q{=}0.97932$};
  \node[vgnode, fill=BuildBlue!45, inner sep=1.8pt] (kf6) at (-1.5,-2.2) {Kf6\\ \scriptsize $n{=}58$\ \ $Q{=}0.96332$};
  \node[vgrefuted, inner sep=1.8pt] (kf5) at (1.5,-2.2) {Kf5\\ \scriptsize $n{=}1$\ \ $Q{=}0.00000$};
  \node[vgrefuted, inner sep=1.8pt] (kd5) at (4.5,-2.2) {Kd5\\ \scriptsize $n{=}1$\ \ $Q{=}0.00000$};
  \vgedgewidth{62}{62}\draw[vgedge, line width=\vgEW pt] (root) -- (kd6);
  \vgedgewidth{58}{62}\draw[vgedge, line width=\vgEW pt] (root) -- (kf6);
  \vgedgewidth{1}{62}\draw[vgedge, line width=\vgEW pt] (root) -- (kf5);
  \vgedgewidth{1}{62}\draw[vgedge, line width=\vgEW pt] (root) -- (kd5);
  \vgdeltadown{kf5}{\scriptsize $d{=}1.000$ refuted}
  \vgdeltadown{kd5}{\scriptsize $d{=}1.000$ refuted}
\end{tikzpicture}
}\\[2ex]
\subcaptionbox{800 nodes: Kd6 N=377, Kf6 N=240, drawing moves stay at N=1}{% GENERATED by tools/make_figures.py — do not edit.
\begin{tikzpicture}
  \node[vgroot] (root) at (0,0) {Root\\ \scriptsize Budget: 800 nodes};
  \node[vgnode, fill=BuildBlue!60, inner sep=1.8pt] (kd6) at (-4.5,-2.2) {Kd6\\ \scriptsize $n{=}377$\ \ $Q{=}0.95610$};
  \node[vgnode, fill=BuildBlue!50, inner sep=1.8pt] (kf6) at (-1.5,-2.2) {Kf6\\ \scriptsize $n{=}240$\ \ $Q{=}0.92742$};
  \node[vgrefuted, inner sep=1.8pt] (kf5) at (1.5,-2.2) {Kf5\\ \scriptsize $n{=}1$\ \ $Q{=}0.00000$};
  \node[vgrefuted, inner sep=1.8pt] (kd5) at (4.5,-2.2) {Kd5\\ \scriptsize $n{=}1$\ \ $Q{=}0.00000$};
  \vgedgewidth{377}{377}\draw[vgedge, line width=\vgEW pt] (root) -- (kd6);
  \vgedgewidth{240}{377}\draw[vgedge, line width=\vgEW pt] (root) -- (kf6);
  \vgedgewidth{1}{377}\draw[vgedge, line width=\vgEW pt] (root) -- (kf5);
  \vgedgewidth{1}{377}\draw[vgedge, line width=\vgEW pt] (root) -- (kd5);
  \node[font=\scriptsize\bfseries, text=PitfallRed, below=8mm of kf5] {Still $n{=}1$ after 672 more nodes!};
\end{tikzpicture}
}
\captionof{figure}{\textbf{FIG-3.1: The refutation ladder.} Between 64 and 128 nodes, Kf5 and Kd5 each receive exactly one visit, return $Q = 0.000$ (Equation~\eqref{eq:E4}), and are permanently abandoned for the next 672 nodes.}
\label{fig:3_1}
\end{minipage}
\end{center}
\end{realdata}

\section{The Score Decay Race}

\begin{established}[PUCT Score Curves]
\begin{center}
% GENERATED by tools/make_figures.py — do not edit.
\begin{tikzpicture}
  \draw[->, WorkGrey] (0,0) -- (7.5,0) node[right, font=\scriptsize] {$N$ (root visits)};
  \draw[->, WorkGrey] (0,0) -- (0,3.5) node[above, font=\scriptsize] {$S$ score};
  \draw[WorkGrey!30, dash pattern=on 1pt off 1.5pt] (0,2.2) -- (7.5,2.2) node[right, font=\scriptsize] {$S\approx 1.0$};
  \draw[thick, BuildBlue] (0,3.2) .. controls (3,2.8) and (5,2.4) .. (7.5,2.2) node[right, font=\scriptsize] {Kd6 / Kf6};
  \draw[thick, PitfallRed, dashed] (0,1.0) .. controls (3.5,1.5) and (5.0,2.3) .. (5.5,2.5);
  \fill[PitfallRed] (5.5,2.5) circle (2pt) node[above left, font=\scriptsize] {N=64..128 (visit 1)};
  \draw[thick, PitfallRed] (5.5,2.5) -- (5.6,0.2) -- (7.5,0.2) node[right, font=\scriptsize] {Kf5 / Kd5 ($Q=0$)};
  \node[font=\scriptsize\itshape, text=WorkGrey, align=center] at (3.5,-0.6) {Curves derived via PUCT formula from measured priors; markers at measured $N=64, 128$.};
\end{tikzpicture}

\captionof{figure}{\textbf{FIG-3.2: PUCT score race.} Solid markers represent measured engine budgets ($N=64, 128$). Dashed curves represent theoretical $S(N) = Q_{\text{FPU}}(N) + U(N)$ derived via the PUCT formula (Equations~\eqref{eq:E10}--\eqref{eq:E11}) as unvisited $U$ grows with $\sqrt{N}$ until a single visit drops $Q$ to zero.}
\label{fig:3_2}
\end{center}
\end{established}

\section{Contrast: When Nothing Matters}

\begin{realdata}[Dead Drawn Positions]
\begin{center}
% GENERATED by tools/make_figures.py — do not edit.
\begin{tikzpicture}
  \node[vgroot] (root) at (0,0) {Root: Opposition Draw\\ \scriptsize Budget: 800 nodes, all 5 moves draw};
  \node[vgnode, fill=BuildBlue!30] (m1) at (-4.8,-2.2) {Kd4\\ \scriptsize $n{=}171$};
  \node[vgnode, fill=BuildBlue!30] (m2) at (-2.4,-2.2) {Ke3\\ \scriptsize $n{=}148$};
  \node[vgnode, fill=BuildBlue!30] (m3) at (0,-2.2) {Kf4\\ \scriptsize $n{=}172$};
  \node[vgnode, fill=BuildBlue!28] (m4) at (2.4,-2.2) {Kd3\\ \scriptsize $n{=}155$};
  \node[vgnode, fill=BuildBlue!28] (m5) at (4.8,-2.2) {Kf3\\ \scriptsize $n{=}152$};
  \draw[vgedge, line width=1.5pt] (root) -- (m1);
  \draw[vgedge, line width=1.4pt] (root) -- (m2);
  \draw[vgedge, line width=1.5pt] (root) -- (m3);
  \draw[vgedge, line width=1.4pt] (root) -- (m4);
  \draw[vgedge, line width=1.4pt] (root) -- (m5);
  \node[font=\scriptsize\itshape, text=WorkGrey, below=4mm of m3] {Flat visit distribution across all legal moves: position is drawn.};
\end{tikzpicture}

\captionof{figure}{\textbf{FIG-3.3: Opposition Draw Contrast.} In a dead drawn position, visit counts split almost evenly across all 5 legal moves ($N = 171, 148, 172, 155, 152$). A flat tree indicates no tactical differentiation.}
\label{fig:3_3}
\end{center}
\end{realdata}

\section{The Honest Caveat}

\begin{established}[Single-Look vs Deep Refutation]
\begin{center}
% GENERATED by tools/make_figures.py — do not edit.
\begin{tikzpicture}
  \node[draw=NavyBlue, fill=PaleGrey, rounded corners=3pt, inner sep=6pt, text width=55mm, font=\small] (left) at (0,0) {
    \textbf{Refutation in 1 Look}\\[0.5ex]
    Value head catches draw/loss immediately on first visit ($Q=0.000$). Search abandons move.
  };
  \node[draw=NavyBlue, fill=PaleGrey, rounded corners=3pt, inner sep=6pt, text width=55mm, font=\small, right=6mm of left] (right) {
    \textbf{Refutation 5 Plies Deep}\\[0.5ex]
    Value head initial look is optimistic. Search must build subtree before truth emerges.
  };
\end{tikzpicture}

\captionof{figure}{\textbf{FIG-3.4: One-look vs deep refutation.} Rapid refutation relies on the value head recognizing drawn/lost structures on the first visit (Equation~\eqref{eq:E2}). Deep tactical refutations require building multi-ply subtrees.}
\label{fig:3_4}
\end{center}
\end{established}

